\section{SYSTEM SETUP AND API REFERENCE}

\subsection{Runtime Environment}

Table~\ref{tab:runtime} lists the software versions used during development
and testing.

\begin{table}[H]
  \centering
  \caption{Key runtime components (representative versions)}
  \label{tab:runtime}
  \begin{tabularx}{\linewidth}{lX}
    \toprule
    \textbf{Component} & \textbf{Notes} \\
    \midrule
    Node.js 18 LTS & Backend runtime; required by H5P server package. \\
    React 18 + Vite & Frontend runtime and build tool (TypeScript used throughout). \\
    SQLite (sqlite3) & Lightweight relational store used in development; Sequelize ORM. \\
    AI tooling & Optional: OpenAI/Anthropic APIs, Ollama local inference, faster-whisper for transcription. \\
    Utilities & Nodemailer (email), Helmet (security), CORS and rate-limiting middleware. \\
    \bottomrule
  \end{tabularx}
\end{table}

\subsection{Starting the Platform}

\begin{enumerate}
  \item \textbf{Install backend dependencies:}
\begin{lstlisting}[language={}]
cd backend && npm install
\end{lstlisting}

  \item \textbf{Configure environment:} Copy \texttt{backend/.env.example}
  to \texttt{backend/.env} and set \texttt{JWT\_SECRET} and
  \texttt{GROQ\_API\_KEY} (required for AI generation).

  \item \textbf{Start the backend} (runs on port~5001 by default):
\begin{lstlisting}[language={}]
cd backend && npm run dev
\end{lstlisting}

  \item \textbf{Install frontend dependencies:}
\begin{lstlisting}[language={}]
cd frontend && npm install
\end{lstlisting}

  \item \textbf{Start the frontend} (runs on port~3002 by default):
\begin{lstlisting}[language={}]
cd frontend && npm run dev
\end{lstlisting}

  \item \textbf{Open the application} at
  \texttt{http://localhost:3002} in a modern web browser.
\end{enumerate}

\subsection{Key API Endpoints Reference}

\footnotesize
\begin{longtable}{>{\raggedright\arraybackslash}p{1.6cm} >{\raggedright\arraybackslash}p{5.0cm} >{\raggedright\arraybackslash}p{8.1cm}}
  \caption{Key API endpoints} \label{tab:api_ref} \\
  \toprule
  \textbf{Method} & \textbf{Path} & \textbf{Description} \\
  \midrule
  \endfirsthead
  \multicolumn{3}{c}{\tablename\ \thetable{} (continued)} \\
  \toprule
  \textbf{Method} & \textbf{Path} & \textbf{Description} \\
  \midrule
  \endhead
  \bottomrule
  \endlastfoot
  POST & /api/auth/register        & Create a new user account. \\
  POST & /api/auth/login            & Authenticate and receive JWT. \\
  GET  & /api/videos                & List authenticated user's videos. \\
  POST & /api/videos                & Upload a new video file. \\
  POST & /api/videos/youtube        & Import a YouTube video by URL. \\
  GET  & /api/videos/:id            & Get a single video's metadata. \\
  DELETE & /api/videos/:id          & Soft-delete a video. \\
  GET  & /api/h5p/:videoId/content  & Get H5P interactions for a video. \\
  POST & /api/h5p/:videoId/content  & Add an H5P interaction manually. \\
  POST & /api/h5p/video/:videoId/export & Export H5P package (binary download). \\
  GET  & /api/ai/status             & Check AI provider availability. \\
  POST & /api/ai/analyze            & Analyse transcript (non-streaming). \\
  POST & /api/ai/analyze-stream     & Analyse transcript (SSE streaming). \\
  POST & /api/ai/inject             & Inject accepted suggestions into H5P. \\
  POST & /api/transcript/youtube    & Extract captions from YouTube video. \\
  GET  & /api/admin/users           & List all users (admin only). \\
  PATCH & /api/admin/users/:id      & Update user role or active status (admin only). \\
  POST & /api/admin/users           & Create a new account (admin only). \\
  GET  & /api/admin/audit-log       & Retrieve admin audit log (admin only). \\
  GET  & /api/admin/settings        & Get all system settings (admin only). \\
  PUT  & /api/admin/settings/:key   & Update a system setting (admin only). \\
  PUT  & /api/users/profile         & Update authenticated user's profile. \\
  POST & /api/auth/verify-email     & Submit 6-digit email verification code. \\
  POST & /api/auth/forgot-password  & Request password reset email. \\
  POST & /api/auth/reset-password   & Submit new password with reset token. \\
  GET  & /api/lti/launch            & LTI~1.3 launch endpoint for LMS integration. \\
\end{longtable}
